\documentclass[12pt,letterpaper]{article}

\usepackage[margin=0.75in]{geometry}
\usepackage{amsmath,amssymb}
\usepackage{enumitem}
\usepackage{environ}
\usepackage{tikz}

\pagestyle{plain}
\setlength{\parindent}{0pt}
\setlength{\parskip}{0.5em}
\setlength{\emergencystretch}{2em}

% Solutions are hidden unless \showsolutions is defined; midterm1_key.tex does
% this and then inputs this file, producing the answer key as a separate PDF.
\newif\ifsolutions
\ifdefined\showsolutions\solutionstrue\fi

\NewEnviron{solution}{%
  \ifsolutions
    \par{\color{blue!55!black}\textbf{Solution.} \BODY\par}%
  \fi}

\newcommand{\points}[1]{(#1 points)}
\newcommand{\workspace}{\vfill\newpage}

\begin{document}

\thispagestyle{empty}

\begin{center}
  {\large\bfseries MATH 15200: Calculus II}\\
  {\large\bfseries Midterm 1}\\[0.3em]
  Tuesday, October 27, 2026
\end{center}

\vspace{0.25em}
\hrule
\vspace{1.5em}

\textbf{Name:} \rule{4in}{0.4pt}

\ifsolutions
  \begin{center}
    {\color{blue!55!black}\large\bfseries ANSWER KEY}
  \end{center}
\fi

\vspace{0.8em}

\textbf{Instructions.} You have 80 minutes. No notes, books, calculators,
phones, or other electronic devices are allowed. Show all of your work and
justify your answers; a correct answer without supporting work may not receive
full credit. Give exact answers (such as $\tfrac{\sqrt{2}}{3}$ or $5\pi$), not
decimal approximations. The last page is extra space. If you run out of room
on a problem, you may continue your answer there; if you do, write on that
problem's page that your answer continues on the last page, and label your
work on the last page with the problem number.

\vfill

\begin{center}
  \renewcommand{\arraystretch}{1.5}
  \begin{tabular}{|c|c|p{1in}|}
    \hline
    \textbf{Problem} & \textbf{Points} & \textbf{Score} \\
    \hline
    1 & 20 & \\ \hline
    2 & 15 & \\ \hline
    3 & 15 & \\ \hline
    4 & 10 & \\ \hline
    5 & 15 & \\ \hline
    6 & 25 & \\ \hline
    \textbf{Total} & 100 & \\ \hline
  \end{tabular}
\end{center}

\vfill
\newpage

\begin{enumerate}[label=\textbf{\arabic*.},leftmargin=2em]
  \item \points{20}
    \begin{enumerate}[label=(\alph*)]
      \item \points{10} Prove directly from the $\varepsilon$--$\delta$
        definition that
        \[
          \lim_{x\to 3}\bigl(x^2-2x\bigr)=3.
        \]

        \begin{solution}
          Let $\varepsilon>0$ and set $\delta=\min\{1,\varepsilon/5\}$.
          Suppose $0<|x-3|<\delta$. Since $|x-3|<1$, we have $2<x<4$, so
          $3<x+1<5$ and $|x+1|<5$. Therefore
          \[
            \bigl|(x^2-2x)-3\bigr|=|x-3|\,|x+1|\le 5|x-3|<5\delta\le\varepsilon.
          \]
        \end{solution}
        \vfill

      \item \points{10} Define a function $f$ by
        \[
          f(x)=
          \begin{cases}
            3, & x\neq 1,\\
            5, & x=1.
          \end{cases}
        \]
        Prove directly from the definition that $f$ is not continuous at $1$.

        \begin{solution}
          Since $f(1)=5$, we must show that $\lim_{x\to 1}f(x)=5$ is false:
          that there is an $\varepsilon>0$ such that for every $\delta>0$,
          some $x$ with $0<|x-1|<\delta$ has $|f(x)-5|\ge\varepsilon$. Take
          $\varepsilon=1$. Given $\delta>0$, let $x=1+\delta/2$. Then
          $0<|x-1|<\delta$ and $x\neq 1$, so
          $|f(x)-5|=|3-5|=2\ge\varepsilon$. Hence $f$ is not continuous
          at $1$. (Also valid, if uniqueness of limits was proved in class:
          $\lim_{x\to 1}f(x)=3$, since any $\delta$ works, and $3\neq f(1)$.)
        \end{solution}
        \vfill
    \end{enumerate}
    \newpage

  \item \points{15}
    Let $f(x)=1/x$ on $[1,2]$, and let $P=\{1,\tfrac43,\tfrac53,2\}$.
    \begin{enumerate}[label=(\alph*)]
      \item \points{10} Compute the lower sum $L_f(P)$ and the upper sum
        $U_f(P)$.

        \begin{solution}
          Each subinterval has length $\tfrac13$. Since $f$ is decreasing on
          $[1,2]$, its minimum on each subinterval is at the right endpoint
          and its maximum is at the left endpoint. Hence
          \[
            L_f(P)=\frac13\Bigl(\frac34+\frac35+\frac12\Bigr)
              =\frac13\cdot\frac{37}{20}=\frac{37}{60},
            \qquad
            U_f(P)=\frac13\Bigl(1+\frac34+\frac35\Bigr)
              =\frac13\cdot\frac{47}{20}=\frac{47}{60}.
          \]
        \end{solution}
        \vfill

      \item \points{5} Use part (a) to find numbers $A$ and $B$ such that
        \[
          A\le \ln(2) \le B,
        \]
        and explain why these inequalities hold.

        \begin{solution}
          Take $A=L_f(P)$ and $B=U_f(P)$. By the definition of the natural
          logarithm, $\ln 2=\int_1^2\frac{dx}{x}$ (or by the Fundamental
          Theorem of Calculus, since $\frac{d}{dx}\ln x=\frac1x$), and the
          integral lies between every lower sum and every upper sum, so
          \[
            \frac{37}{60}\le\ln 2=\int_1^2\frac{dx}{x}\le\frac{47}{60}.
          \]
        \end{solution}
        \vfill
    \end{enumerate}
    \newpage

  \item \points{15}
    For $x>0$, let
    \[
      F(x)=\int_{x}^{x^2}\sqrt{t^3+8}\,dt.
    \]
    \begin{enumerate}[label=(\alph*)]
      \item \points{10} Find $F'(x)$.

        \begin{solution}
          By the Fundamental Theorem of Calculus and the chain rule (as in
          Homework~2, Problem~1),
          \[
            F'(x)=\sqrt{(x^2)^3+8}\cdot 2x-\sqrt{x^3+8}\cdot 1
              =2x\sqrt{x^6+8}-\sqrt{x^3+8}.
          \]
        \end{solution}
        \vfill

      \item \points{5} Find an equation of the tangent line to the graph of
        $F$ at $x=1$.

        \begin{solution}
          $F(1)=\int_1^1\sqrt{t^3+8}\,dt=0$ and $F'(1)=2\sqrt9-\sqrt9=3$, so
          the tangent line is $y=3(x-1)$, that is, $y=3x-3$.
        \end{solution}
        \vfill
    \end{enumerate}
    \newpage

  \item \points{10}
    Evaluate the following integrals.
    \begin{enumerate}[label=(\alph*)]
      \item \points{5}\quad
        $\displaystyle\int\cos x\,\sqrt{1+\sin x}\,dx$

        \begin{solution}
          Let $u=1+\sin x$, so $du=\cos x\,dx$. Then
          \[
            \int\cos x\,\sqrt{1+\sin x}\,dx=\int\sqrt{u}\,du
            =\frac23u^{3/2}+C=\frac23(1+\sin x)^{3/2}+C.
          \]
        \end{solution}
        \vfill

      \item \points{5}\quad
        $\displaystyle\int_1^4\frac{dx}{\sqrt{x}\,\bigl(1+\sqrt{x}\bigr)^2}$

        \begin{solution}
          Let $u=1+\sqrt{x}$, so $du=\dfrac{dx}{2\sqrt{x}}$; when $x=1$,
          $u=2$, and when $x=4$, $u=3$. Then
          \[
            \int_1^4\frac{dx}{\sqrt{x}\,\bigl(1+\sqrt{x}\bigr)^2}
            =\int_2^3\frac{2\,du}{u^2}
            =\Bigl[-\frac{2}{u}\Bigr]_2^3=-\frac23+1=\frac13.
          \]
        \end{solution}
        \vfill
    \end{enumerate}
    \newpage

  \item \points{15}
    Find the total (unsigned) area enclosed by the curves $y=x^3$ and
    $y=x^2+2x$.

    \begin{solution}
      The curves meet where $x^3=x^2+2x$, that is, $x(x+1)(x-2)=0$, so at
      $x=-1$, $0$, and $2$. On $[-1,0]$ the cubic is on top (at
      $x=-\tfrac12$: $-\tfrac18>-\tfrac34$), and on $[0,2]$ the parabola is
      on top (at $x=1$: $3>1$). Hence the area is
      \begin{align*}
        A&=\int_{-1}^{0}\bigl(x^3-x^2-2x\bigr)\,dx
           +\int_0^2\bigl(x^2+2x-x^3\bigr)\,dx\\
         &=\Bigl[\frac{x^4}{4}-\frac{x^3}{3}-x^2\Bigr]_{-1}^{0}
           +\Bigl[\frac{x^3}{3}+x^2-\frac{x^4}{4}\Bigr]_0^2
          =\frac{5}{12}+\frac83=\frac{37}{12}.
      \end{align*}
    \end{solution}
    \workspace

  \item \points{25}
    Let $r>0$ and $h>0$. Revolving the disk
    $x^2+\bigl(y-(h+r)\bigr)^2\le r^2$, which has radius $r$ and center
    $(0,h+r)$, about the $x$-axis produces a solid called a \emph{torus};
    the hole in the middle of the torus has radius $h$ (see the figure).
    Find the volume of the torus in terms of $r$ and $h$.

    \begin{center}
      % Tube radius r, hole radius h, and viewing elevation in degrees.
      \pgfmathsetmacro{\tubeR}{0.65}
      \pgfmathsetmacro{\holeR}{1.3}
      \pgfmathsetmacro{\elev}{30}
      \pgfmathsetmacro{\ctrR}{\holeR+\tubeR}
      \begin{tikzpicture}[
          axis/.style={thin, dash dot},
          declare function={
            % Silhouette angle around the tube at angle t around the axis.
            sil(\t) = atan(sin(\t)*cot(\elev));
            outx(\t) = (\ctrR + \tubeR*cos(sil(\t)))*cos(\t);
            outy(\t) = (\ctrR + \tubeR*cos(sil(\t)))*sin(\t)*sin(\elev)
                       + \tubeR*sin(sil(\t))*cos(\elev);
            inx(\t) = (\ctrR - \tubeR*cos(sil(\t)))*cos(\t);
            iny(\t) = (\ctrR - \tubeR*cos(sil(\t)))*sin(\t)*sin(\elev)
                      - \tubeR*sin(sil(\t))*cos(\elev);
          },
        ]
        % The torus is drawn with its axis vertical, seen from \elev degrees
        % above its middle plane, and then turned a quarter turn clockwise so
        % that its axis is horizontal.
        \begin{scope}[rotate=-90]
          % The angle ranges below were computed for h = 2r and a 30 degree
          % elevation: the inner rim is visible for 195.5 <= t <= 344.5
          % (front) and 30.6 <= t <= 149.4 (back), and the two cross at
          % t = 210.6, 329.4 (front) and t = 30.6, 149.4 (back).
          \fill[gray!15] plot[domain=0:360, samples=181, variable=\t]
            ({outx(\t)}, {outy(\t)}) -- cycle;
          \fill[white] plot[domain=210.6:329.4, samples=60, variable=\t]
              ({inx(\t)}, {iny(\t)})
            -- plot[domain=30.6:149.4, samples=60, variable=\t]
              ({inx(\t)}, {iny(\t)}) -- cycle;
          \draw plot[domain=0:360, samples=181, variable=\t]
            ({outx(\t)}, {outy(\t)}) -- cycle;
          \draw plot[domain=195.5:344.5, samples=80, variable=\t]
            ({inx(\t)}, {iny(\t)});
          \draw plot[domain=30.6:149.4, samples=60, variable=\t]
            ({inx(\t)}, {iny(\t)});
          % Axis of revolution, hidden where it passes behind the front of
          % the torus.
          \draw[axis] (0,{\ctrR*sin(\elev)+\tubeR+0.6})
            -- (0,{\tubeR-\ctrR*sin(\elev)});
          \draw[axis] (0,{-\ctrR*sin(\elev)-\tubeR})
            -- (0,{-\ctrR*sin(\elev)-\tubeR-0.5});
          \draw[->, thin] (0,{\ctrR*sin(\elev)+\tubeR+0.35})
            ++(200:0.35 and 0.12)
            arc[start angle=200, end angle=500, x radius=0.35, y radius=0.12];
        \end{scope}
        \node[font=\small] at (0,{-\ctrR-\tubeR-0.45}) {the torus};

        % The disk being revolved, in the xy-plane; the axis of revolution is
        % the x-axis.
        \begin{scope}[xshift=4.7cm, yshift=-1.5cm]
          \fill[gray!15] (0,\ctrR) circle[radius=\tubeR];
          \draw (0,\ctrR) circle[radius=\tubeR];
          \draw[->] (-1.4,0) -- (1.6,0) node[right] {$x$};
          \draw[->] (0,-0.4) -- (0,{\ctrR+\tubeR+0.45}) node[above] {$y$};
          \draw[->, thin] (1.2,0) ++(110:0.12 and 0.35)
            arc[start angle=110, end angle=410, x radius=0.12, y radius=0.35];
          \fill (0,\ctrR) circle[radius=1.2pt];
          \draw (0,\ctrR) -- node[above, inner sep=1pt] {$r$} ++(\tubeR,0);
          \draw[thin] (0,\holeR) -- (-0.55,\holeR);
          \draw[<->] (-0.4,0) -- node[left] {$h$} (-0.4,\holeR);
        \end{scope}
        \node[font=\small] at (4.85,{-\ctrR-\tubeR-0.45})
          {the disk and the axis};
      \end{tikzpicture}
    \end{center}

    \emph{Hint:} Use the washer method, and note that
    $\int_{-r}^{r}\sqrt{r^2-x^2}\,dx$ is the area of a semicircle of
    radius~$r$.

    \begin{solution}
      The disk is bounded above and below by the semicircles
      $y=(h+r)\pm\sqrt{r^2-x^2}$, for $-r\le x\le r$. So the cross-section
      of the torus perpendicular to the $x$-axis at $x$ is a washer with
      outer radius $(h+r)+\sqrt{r^2-x^2}$ and inner radius
      $(h+r)-\sqrt{r^2-x^2}$, and its area is
      \[
        A(x)=\pi\Bigl[\bigl(h+r+\sqrt{r^2-x^2}\bigr)^2
          -\bigl(h+r-\sqrt{r^2-x^2}\bigr)^2\Bigr]
        =4\pi(h+r)\sqrt{r^2-x^2}.
      \]
      The integral $\int_{-r}^{r}\sqrt{r^2-x^2}\,dx$ is the area of a
      semicircle of radius $r$, namely $\pi r^2/2$. Hence
      \[
        V=\int_{-r}^{r}4\pi(h+r)\sqrt{r^2-x^2}\,dx
         =4\pi(h+r)\cdot\frac{\pi r^2}{2}=2\pi^2r^2(h+r).
      \]
      (Shells also work:
      $V=\int_h^{h+2r}2\pi y\cdot 2\sqrt{r^2-(y-h-r)^2}\,dy$; after
      substituting $u=y-(h+r)$, the term $\int_{-r}^{r}u\sqrt{r^2-u^2}\,du$
      vanishes by symmetry.) As a check, $r=1$ and $h=2$ give $6\pi^2$, the
      answer to Homework~3, Problem~2.
    \end{solution}
    \workspace
\end{enumerate}

\begin{center}
  {\Large\bfseries Extra Space}\\[0.3em]
  \noindent\rule{\linewidth}{0.4pt}
\end{center}

% \framebox[\textwidth]{\rule{0pt}{8.65in}}

\end{document}
